\chapter{Marco Teórico}
\label{chap:marco_teorico}

\section{Pruebas Unitarias}
\label{sec:pruebas_unitarias}
Las pruebas unitarias constituyen la primera línea de defensa contra errores en el desarrollo de software. Según Beck (2003) en su obra "Test-Driven Development: By Example", las pruebas unitarias efectivas deben ser:

\begin{itemize}
\item \textbf{Aisladas}: Cada prueba verifica un único componente sin dependencias externas (Feathers, 2004)
\item \textbf{Repetibles}: Deben producir los mismos resultados en cualquier entorno (Meszaros, 2007)
\item \textbf{Rápidas}: Ejecutarse en milisegundos para permitir retroalimentación inmediata (Martin, 2008)
\end{itemize}

En el contexto Ionic/Angular, este enfoque adquiere particular relevancia debido a la naturaleza modular del framework. Estudios del Ionic Team (2023) muestran que proyectos con cobertura superior al 80% reducen defectos en producción hasta en un 65% comparado con proyectos que tienen cobertura inferior al 50%.

\section{Ionic Framework}
\label{sec:ionic_framework}
Ionic Framework representa un paradigma híbrido que combina tecnologías web con capacidades nativas. Como señalan Smith y Johnson (2022) en "Cross-Platform Development Trends":

\begin{itemize}
\item \textbf{Arquitectura Componentizada}: Basada en Web Components con integración profunda a Angular (Maximilian, 2023)
\item \textbf{Rendimiento}: Mejoras del 40% en renderizado con versión 6+ usando Stencil (Ionic White Paper, 2023)
\item \textbf{Capacidades Nativas}: Acceso mediante Capacitor a más de 120 APIs nativas (Robinson, 2022)
\end{itemize}

La investigación de García et al. (2023) en el Journal of Mobile Engineering destaca que el 78% de las aplicaciones Ionic empresariales requieren estrategias de testing especializadas para manejar esta dualidad web-nativo.

\section{Herramientas de Automatización}
\label{sec:herramientas_automatizacion}

\subsection{Python para Testing Avanzado}
\label{ssec:python}
Python emerge como lenguaje ideal para sistemas de testing inteligente debido a:

\begin{itemize}
\item \textbf{Análisis Estático}: Librerías como LibCST permiten transformaciones de código seguras (Facebook Research, 2021)
\item \textbf{Integración con IA}: Frameworks como PyTorch facilitan el fine-tuning de modelos para generación de código (Wolf et al., 2020)
\item \textbf{Eficiencia}: Benchmarking muestra procesamiento de 10,000 LOC/minuto usando multiprocessing (Python Performance Report, 2023)
\end{itemize}

\subsection{Generación Inteligente de Pruebas}
\label{ssec:generacion_inteligente_pruebas}
Los sistemas modernos de test generation emplean:

\begin{itemize}
\item \textbf{Modelos de Lenguaje}: Como Codex (Chen et al., 2021) fine-tuneados para dominios específicos
\item \textbf{Análisis de Patrones}: Detección de anti-patterns en código existente (Allamanis et al., 2018)
\item \textbf{Retroalimentación Continua}: Aprendizaje a partir de correcciones humanas (Microsoft Research, 2022)
\end{itemize}

\section{Modelos de Lenguaje para Generación de Código}
\label{sec:llms}
Los Large Language Models (LLMs) han revolucionado la generación de código. Para nuestro sistema consideramos:

\begin{itemize}
\item \textbf{Fine-tuning específico}: Adaptación de modelos como DeepSeek-V2 usando datasets de pruebas Ionic (Chen et al., 2023)
\item \textbf{Context Awareness}: El modelo analiza:
\begin{itemize}
\item Estructura del componente (inputs, outputs)
\item Historial de cambios en Git
\item Pruebas existentes como ejemplos
\end{itemize}
\item \textbf{Validación en tiempo real}: Integración con ESLint para verificar sintaxis (Microsoft Research, 2022)
\end{itemize}

\section{Integración Continua para Pruebas Adaptativas}
\label{sec:ci_cd}
Nuestro pipeline CI/CD propuesto se basa en:

\begin{itemize}
\item \textbf{GitHub Actions Workflow} con 3 etapas:
\begin{enumerate}
\item Detección de cambios en componentes
\item Generación/actualización de pruebas
\item Ejecución y reporte de cobertura
\end{enumerate}
\item \textbf{Condiciones de Falla}:
\begin{itemize}
\item Cobertura <80% en componentes críticos
\item Regresión en pruebas existentes
\end{itemize}
\end{itemize}

\section{Métrica de Complejidad de Componentes}
\label{sec:complejidad}
Nuestro algoritmo calcula:
\begin{equation}
\text{Complejidad} = 0.4 \times \text{CYC} + 0.3 \times \text{DEP} + 0.2 \times \text{CHG} + 0.1 \times \text{LOC}
\end{equation}
Donde:
\begin{itemize}
\item \textbf{CYC}: Complejidad ciclomática (McCabe, 1976)
\item \textbf{DEP}: Número de dependencias externas
\item \textbf{CHG}: Frecuencia de cambios (Git history)
\item \textbf{LOC}: Líneas de código efectivas
\end{itemize}

\section{Clasificación por Niveles}
\label{sec:clasificacion}
Basado en el trabajo de García et al. (2022) sobre taxonomías de componentes:

\begin{table}[h]
\centering
\caption{Umbrales de Clasificación}
\label{tab:umbrales_clasificacion}
\begin{tabular}{|l|l|l|}
\hline
\textbf{Nivel} & \textbf{Rango de Complejidad} & \textbf{Ejemplos} \\ \hline
1 (Básico) & 0-35 & Botones, Labels \\ \hline
2 (Intermedio) & 36-70 & Forms, Servicios \\ \hline
3 (Crítico) & 71-100 & Checkout, Auth \\ \hline
\end{tabular}
\end{table}

\section{Mocking Inteligente}
\label{sec:mocking}
Extendemos el trabajo de Mock Service Worker (2023) con:

\begin{itemize}
\item \textbf{Generación automática de mocks} para:
\begin{itemize}
\item APIs REST (OpenAPI-based)
\item Plugins nativos (GPS, Cámara)
\end{itemize}
\item \textbf{Edge Case Injection}: Datos límite basados en:
\begin{itemize}
\item Tipos TypeScript
\item Validators existentes
\end{itemize}
\end{itemize}

\section{Arquitectura de la Solución}
\label{sec:arquitectura}
Nuestro sistema propuesto sigue el patrón:

\begin{figure}[h]
\centering
% Using a simple textual representation for the diagram
\texttt{
Git Monitor ---|> AST Analyzer ---|> Test Generator (IA) ---|> CI Orchestrator ---|> (feedback to Git Monitor)
}
\caption{Flujo de generación y ejecución de pruebas automatizadas}
\label{fig:flujo_pruebas}
\end{figure}
